% !PRINTED 613
\weberpage{613}

If we now form, according to (4), the composition
\[
\lambda \lambda_0^{-h} = \left(z,\, z + a_0^{-h} b + \frac{a_0^{-h} - 1}{a_0 - 1} b_0\right),
\]
then we obtain the identical substitution when, and only when,
\[
\frac{b_0}{a_0 - 1} \equiv \frac{b}{a_0^h - 1}\ (\mathrm{mod}\ n).
\]

If this congruence is satisfied for \gesperrt{all} substitutions $\lambda$ of $L$, then the element $\dfrac{b_0}{a_0 - 1}$ remains unchanged by the whole group, and $L$ is intransitive. Consequently, in a transitive group $L$ there is always an element $\lambda$ for which $\lambda \lambda_0^{-h}$ is not the identical substitution, although the index is zero, and $L$ therefore contains the whole cyclic group.

If one forms the various powers of $\lambda_0$, it turns out that the exponent $h$ in the substitutions $\lambda$ of the group $L$ can take every value. If $e$ is the least positive exponent that satisfies the condition $a_0^e \equiv 1\ (\mathrm{mod}\ n)$, then the values $h = 0, 1, 2 \ldots e - 1$ therefore occur, and here $e$ is a divisor of $n - 1$. From the composed substitution
\[
(z, a_0^h z + b)(z, z + 1) = (z, a_0^h z + b + 1),
\]
which after all must also occur in $L$, one sees further that with every value of $h$ each of the values $b = 0, 1, 2, \ldots n - 1$ occurs in combination, and the whole group $L$ therefore consists of the substitutions
\weq{5}{\lambda = (z, a_0^h z + b), \quad \begin{array}{l} h = 0, 1 \ldots e - 1 \\ b = 0, 1 \ldots n - 1' \end{array}}
% [sic] source prints a raised apostrophe-like mark after "n-1" in the b-row
and is of degree $en$. Conversely, every system of substitutions of this form constitutes a group.

Of the linear groups the following theorem now holds:

I. \gesperrt{If a transitive linear group $L$ is a normal divisor of another permutation group $P$ of the same $n$ figures, then $P$ is also a linear group.}

The theorem can be proved as follows:
If $\pi$ is an arbitrary permutation:
\[
\pi = \begin{pmatrix} 0, & 1, & 2 \ldots n - 1 \\ a_0, & a_1, & a_2 \ldots a_{n-1} \end{pmatrix},
\]
%CONT>>
